% Source: Gerzensee "Calculus" slides 65-98 (the framing question about the mean
% and variance of a normal; Proposition 3.11 on Riemann integrability;
% Definitions 3.7 and 3.9 and the fundamental theorem of calculus; integration by
% parts, Proposition 3.15, its product-rule derivation, the "integration of f g"
% form and the graphical derivation; the exponential distribution and its two
% moments; the log x exercise; integration by substitution, Proposition 3.16, the
% (2t+3)^4 and x exp(-x^2) exercises; the mean and variance of the normal; the
% Gaussian integral by polar coordinates with the Jacobian determinant; the
% Leibniz formula, Theorem 3.3, Example 3.22 on continuation values, and
% Exercise 3.6).
% Supplementary: Fukuda (2026), Ch. 3 (integration) and Sec. 3.6 (Leibniz).
% External references actually cited in this chapter: Le Gall (2022), Rudin
% (1976).

\chapter{Integration}
\label{ch:integration}

\section{Motivation: Moments of a Distribution}\label{sec:int-motivation}

Expectations are integrals. The question that organises this chapter is the one
posed on \autocite[slide 65]{fukuda2026calculus}: if $X$ is normally distributed
with density
\begin{equation}\label{eq:normal-density}
	f(x)=\frac{1}{\sqrt{2\pi}\,\sigma}\exp\left(-\frac{(x-\mu)^2}{2\sigma^2}\right),
\end{equation}
confirm that $\E[X]=\mu$ and $\Var[X]=\sigma^2$. Doing so requires three tools
--- integration by parts, integration by substitution, and a change of variables
in the plane --- and, before any of them, a definition of the integral precise
enough that the manipulations are licensed rather than merely plausible.

\section{The Riemann Integral}\label{sec:int-riemann}

\begin{definition}[Partitions and Darboux sums]\label{def:darboux}
	Let $f\colon[a,b]\to\R$ be bounded. A \dfn{partition} of $[a,b]$ is a finite set
	$P=\{a=t_0<t_1<\dots<t_k=b\}$. Writing
	$m_i\eqdef\inf_{t\in[t_{i-1},t_i]}f(t)$ and
	$M_i\eqdef\sup_{t\in[t_{i-1},t_i]}f(t)$, the \dfn{lower} and \dfn{upper Darboux
		sums} are
	\[
		L(f,P)\eqdef\sum_{i=1}^{k}m_i(t_i-t_{i-1}),
		\qquad
		U(f,P)\eqdef\sum_{i=1}^{k}M_i(t_i-t_{i-1}).
	\]
	The function is \dfnas{Riemann integrable}{Riemann integrable} on $[a,b]$ if
	$\sup_PL(f,P)=\inf_PU(f,P)$, and the common value is $\int_a^bf(x)\,\diff x$.
	\nidx{integral}{$\int_a^bf$, Riemann integral}
\end{definition}

Boundedness is part of the definition: an unbounded $f$ has $M_i=+\infty$ on some
subinterval of every partition. Integrals over unbounded domains, or of unbounded
functions, are defined as limits in \autoref{def:improper}.

\begin{proposition}[Riemann's criterion]\label{prop:riemann-criterion}
	A bounded $f\colon[a,b]\to\R$ is Riemann integrable if and only if for every
	$\varepsilon>0$ there is a partition $P$ with $U(f,P)-L(f,P)<\varepsilon$.
\end{proposition}

\begin{proof}
	Refining a partition cannot increase $U$ or decrease $L$, and any two partitions
	have a common refinement, so $L(f,P)\leq U(f,Q)$ for all $P,Q$ and hence
	$\sup_PL\leq\inf_QU$. If the criterion holds, the two are within
	$\varepsilon$ of each other for every $\varepsilon$, hence equal. Conversely, if
	they are equal to $I$, choose $P_1$ with $L(f,P_1)>I-\varepsilon/2$ and $P_2$
	with $U(f,P_2)<I+\varepsilon/2$ and take their common refinement.
\end{proof}

\begin{proposition}[Sufficient conditions for integrability]\label{prop:integrable}
	Let $f\colon[a,b]\to\R$.
	\begin{enumerate}[(i)]
		\item If $f$ is continuous, it is Riemann integrable.
		\item If $f$ is bounded with finitely many discontinuities, it is Riemann
		      integrable.
		\item If $f$ is monotone, it is Riemann integrable.
	\end{enumerate}
\end{proposition}

\begin{proof}
	(i) $[a,b]$ is compact, so $f$ is bounded by
	\autoref{cor:weierstrass-continuous} and uniformly continuous by
	\autoref{thm:heine-cantor}. Given $\varepsilon>0$ choose $\delta>0$ with
	$\abs{f(s)-f(t)}<\varepsilon/(b-a)$ whenever $\abs{s-t}<\delta$, and take any
	partition of mesh less than $\delta$. On each subinterval the supremum and
	infimum are attained (\autoref{cor:weierstrass-continuous}) at points less than
	$\delta$ apart, so $M_i-m_i\leq\varepsilon/(b-a)$ and
	\[
		U(f,P)-L(f,P)=\sum_i(M_i-m_i)(t_i-t_{i-1})
		\leq\frac{\varepsilon}{b-a}\sum_i(t_i-t_{i-1})=\varepsilon .
	\]
	Apply \autoref{prop:riemann-criterion}.

	(ii) Let $D$ be the finite set of discontinuities and $\abs{f}\leq M$. Cover
	$D$ by finitely many intervals of total length less than
	$\varepsilon/(4M)$; on those intervals the contribution to $U-L$ is at most
	$2M\cdot\varepsilon/(4M)=\varepsilon/2$. On the complement, a finite union of
	closed intervals on which $f$ is continuous, apply (i) to get a further
	$\varepsilon/2$.

	(iii) Let $f$ be non-decreasing and take the uniform partition of mesh
	$(b-a)/k$. Then $M_i=f(t_i)$ and $m_i=f(t_{i-1})$, so the sum telescopes:
	\[
		U(f,P)-L(f,P)=\frac{b-a}{k}\sum_{i=1}^{k}\bigl(f(t_i)-f(t_{i-1})\bigr)
		=\frac{b-a}{k}\bigl(f(b)-f(a)\bigr)\longrightarrow0 .
	\]
\end{proof}

This is Proposition~3.11 of \textcite{fukuda2026notes}
\autocite[slide 66]{fukuda2026calculus}, with the proofs supplied. Part (iii) is
worth noting: a monotone function may have countably many discontinuities and is
still integrable, so (ii) and (iii) are genuinely different criteria. The exact
characterisation --- integrable if and only if the set of discontinuities has
Lebesgue measure zero --- is Lebesgue's criterion, for which see
\textcite[Thm.~11.33]{rudin1976principles}; it is the point at which the Riemann
theory is superseded by the Lebesgue theory of \autoref{ch:probability}.

\begin{proposition}[Basic properties]\label{prop:integral-properties}
	Let $f,g$ be Riemann integrable on $[a,b]$ and $\alpha,\beta\in\R$. Then
	\begin{enumerate}[(i)]
		\item $\alpha f+\beta g$ is integrable with
		      $\int_a^b(\alpha f+\beta g)=\alpha\int_a^bf+\beta\int_a^bg$;
		\item for $c\in(a,b)$, $\int_a^bf=\int_a^cf+\int_c^bf$;
		\item $f\leq g$ implies $\int_a^bf\leq\int_a^bg$; in particular
		      $\bigl|\int_a^bf\bigr|\leq\int_a^b\abs{f}$;
		\item if $f$ is continuous there is $\xi\in[a,b]$ with
		      $\int_a^bf=f(\xi)(b-a)$.
	\end{enumerate}
\end{proposition}

\begin{proof}
	(i)--(iii) follow from the corresponding statements for Darboux sums and
	\autoref{prop:riemann-criterion}; see
	\textcite[Thm.~6.12]{rudin1976principles}. For (iv), let
	$m=\min_{[a,b]}f$ and $M=\max_{[a,b]}f$, both attained by
	\autoref{cor:weierstrass-continuous}. By (iii),
	$m(b-a)\leq\int_a^bf\leq M(b-a)$, so $\frac{1}{b-a}\int_a^bf\in[m,M]$, and the
	intermediate value theorem (\autoref{cor:ivt}) supplies $\xi$ with
	$f(\xi)$ equal to it.
\end{proof}

\section{The Fundamental Theorem of Calculus}\label{sec:int-ftc}

\begin{definition}[Indefinite integral and antiderivative]\label{def:antiderivative}
	Let $f\colon[a,b]\to\R$ be integrable. The \dfn{indefinite integral} of $f$ is
	$F(x)\eqdef\int_a^xf(t)\,\diff t$ for $x\in[a,b]$. A function
	$G\colon[a,b]\to\R$ is an \dfnas{antiderivative}{antiderivative} (a primitive)
	of $f$ if $G'(x)=f(x)$ for every $x\in(a,b)$. Any two antiderivatives differ by
	a constant, by \autoref{thm:mvt}, which is why the generic antiderivative is
	written $\int f(x)\,\diff x=G(x)+c$.
\end{definition}

\begin{theorem}[Fundamental theorem of calculus]\label{thm:ftc}
	Let $f\colon[a,b]\to\R$ be Riemann integrable and $F(x)=\int_a^xf(t)\,\diff t$.
	\begin{enumerate}[(i)]
		\item $F$ is Lipschitz, hence continuous, on $[a,b]$; and $F$ is
		      differentiable with $F'(x_0)=f(x_0)$ at every point $x_0$ at which $f$ is
		      continuous. In particular, if $f$ is continuous on $[a,b]$ then $F'=f$
		      throughout.
		\item If $G$ is any antiderivative of $f$ on $[a,b]$ and $f$ is integrable,
		      then
		      \[
			      \int_a^bf(t)\,\diff t=G(b)-G(a)\eqdef\bigl[G(t)\bigr]_a^b .
		      \]
	\end{enumerate}
\end{theorem}

\begin{proof}
	(i) With $\abs{f}\leq M$, \autoref{prop:integral-properties}(ii)--(iii) gives
	$\abs{F(y)-F(x)}=\bigl|\int_x^yf\bigr|\leq M\abs{y-x}$, which is the Lipschitz
	bound. Let $f$ be continuous at $x_0$ and $\varepsilon>0$; choose $\delta>0$
	with $\abs{f(t)-f(x_0)}<\varepsilon$ for $\abs{t-x_0}<\delta$. For
	$0<\abs{h}<\delta$,
	\[
		\left|\frac{F(x_0+h)-F(x_0)}{h}-f(x_0)\right|
		=\left|\frac{1}{h}\int_{x_0}^{x_0+h}\bigl(f(t)-f(x_0)\bigr)\,\diff t\right|
		\leq\varepsilon,
	\]
	using \autoref{prop:integral-properties}(iii) and
	$\frac1h\int_{x_0}^{x_0+h}\,\diff t=1$.

	(ii) Let $P=\{a=t_0<\dots<t_k=b\}$ be any partition. By
	\autoref{thm:mvt} applied to $G$ on each $[t_{i-1},t_i]$ there is
	$\xi_i\in(t_{i-1},t_i)$ with
	$G(t_i)-G(t_{i-1})=G'(\xi_i)(t_i-t_{i-1})=f(\xi_i)(t_i-t_{i-1})$. Summing, the
	left side telescopes to $G(b)-G(a)$, and since $m_i\leq f(\xi_i)\leq M_i$,
	\[
		L(f,P)\leq G(b)-G(a)\leq U(f,P)
	\]
	for every $P$. Integrability makes the two outer quantities have the common
	supremum and infimum $\int_a^bf$, so $G(b)-G(a)=\int_a^bf$.
\end{proof}

The two halves say different things, and the slides' formulation
$\deriv{}{x}\int f(x)\,\diff x=f$ and $\int f'(x)\,\diff x=f+c$
\autocite[slide 68]{fukuda2026calculus} compresses them. Part (i) constructs an
antiderivative and requires continuity of $f$ at the point in question; part (ii)
evaluates an integral given an antiderivative and requires only integrability of
$f$ together with the existence of $G$. Neither implies the other: there are
differentiable $G$ whose derivative is not Riemann integrable, for which (ii)
does not apply, and integrable $f$ with no antiderivative, for which (i) fails at
the discontinuities.

\subsection{Integration by parts}

\begin{theorem}[Integration by parts]\label{thm:by-parts}
	Let $u,v\colon[a,b]\to\R$ be differentiable with $u'$ and $v'$ Riemann
	integrable. Then
	\begin{equation}\label{eq:by-parts}
		\int_a^bu'(x)v(x)\,\diff x=\bigl[u(x)v(x)\bigr]_a^b-\int_a^bu(x)v'(x)\,\diff x .
	\end{equation}
\end{theorem}

\begin{proof}
	By the product rule (\autoref{prop:deriv-product}), $uv$ is differentiable with
	$(uv)'=u'v+uv'$, and the right side is integrable as a sum of products of
	integrable functions. Apply \autoref{thm:ftc}(ii) to $uv$ and then
	\autoref{prop:integral-properties}(i).
\end{proof}

Identity \eqref{eq:by-parts} is usually applied in the form of
\autocite[slide 71]{fukuda2026calculus}: given a product $fg$ to integrate, take
an antiderivative $F$ of $f$ and differentiate $g$,
\[
	\int_a^bf(x)g(x)\,\diff x=\bigl[F(x)g(x)\bigr]_a^b-\int_a^bF(x)g'(x)\,\diff x ,
\]
choosing the split so that $Fg'$ is simpler than $fg$. Two standard choices:
$g=\log x$ with $f=1$ reduces $\int\log x\,\diff x$ to $\int1\,\diff x$, giving
$x(\log x-1)+c$ \autocite[slide 78]{fukuda2026calculus}; and $g$ a power with
$f$ an exponential reduces the power by one at each application, which is the
recursion used in \autoref{sec:int-exponential}.

\subsection{Substitution}

\begin{theorem}[Change of variable in one dimension]\label{thm:substitution}
	Let $\varphi\colon[a,b]\to\R$ be continuously differentiable and let $f$ be
	continuous on an interval containing $\varphi([a,b])$. Then
	\begin{equation}\label{eq:substitution}
		\int_{\varphi(a)}^{\varphi(b)}f(x)\,\diff x=\int_a^bf\bigl(\varphi(t)\bigr)\varphi'(t)\,\diff t .
	\end{equation}
\end{theorem}

\begin{proof}
	Let $F$ be an antiderivative of $f$, which exists by \autoref{thm:ftc}(i) since
	$f$ is continuous. By the chain rule (\autoref{prop:chain-rule-1d}),
	$(F\circ\varphi)'(t)=F'(\varphi(t))\varphi'(t)=f(\varphi(t))\varphi'(t)$, which
	is continuous hence integrable. Applying \autoref{thm:ftc}(ii) twice,
	\[
		\int_a^bf\bigl(\varphi(t)\bigr)\varphi'(t)\,\diff t
		=\bigl[F(\varphi(t))\bigr]_a^b=F(\varphi(b))-F(\varphi(a))
		=\int_{\varphi(a)}^{\varphi(b)}f(x)\,\diff x .
	\]
\end{proof}

This is Proposition~3.16 of \textcite{fukuda2026notes}
\autocite[slide 82]{fukuda2026calculus}. Note what the theorem does \emph{not}
require: $\varphi$ need not be injective, and the limits of integration are
$\varphi(a)$ and $\varphi(b)$ in that order, which may be reversed relative to
$a<b$. The mnemonic $x=\varphi(t)$, $\diff x=\varphi'(t)\,\diff t$
\autocite[slide 81]{fukuda2026calculus} is the identity
\eqref{eq:substitution} written in differentials; the substance is the chain
rule, and the differentials are notation for it.

\begin{eg}[Two substitutions]\label{eg:substitutions}
	With $\varphi(t)=2t+3$, $\varphi'=2$, and $f(x)=x^4$,
	\[
		\int_1^2(2t+3)^4\,\diff t=\frac12\int_5^7x^4\,\diff x
		=\frac{1}{10}\bigl[x^5\bigr]_5^7=\frac{7^5-5^5}{10}=\frac{6841}{10} .
	\]
	With $\varphi(x)=-x^2$, $\varphi'(x)=-2x$, and $f(y)=e^y$,
	\[
		\int xe^{-x^2}\,\diff x=-\frac12\int e^{y}\,\diff y=-\frac12e^{-x^2}+c .
	\]
	These are the computations of \autocite[slides 83--84]{fukuda2026calculus}. The
	second is the pattern used for the normal density in
	\autoref{sec:int-normal}: an integrand of the form $g'(x)h(g(x))$ integrates to
	$H(g(x))$ with $H'=h$.
\end{eg}

\subsection{Improper integrals}

\begin{definition}[Improper integral]\label{def:improper}
	Let $f$ be Riemann integrable on $[a,c]$ for every $c>a$. Then
	\[
		\int_a^{\infty}f(x)\,\diff x\eqdef\lim_{c\to\infty}\int_a^cf(x)\,\diff x
	\]
	when the limit exists in $\R$, and analogously at $-\infty$. The integral over
	$\R$ is defined as $\int_{-\infty}^0f+\int_0^{\infty}f$, requiring
	\emph{both} limits to exist separately. The integral \dfnas{converges
		absolutely}{absolute convergence} if $\int\abs{f}<\infty$, and absolute
	convergence implies convergence.
\end{definition}

The requirement that the two halves converge separately is not pedantry. For
$f(x)=x$ the symmetric limit $\lim_{c\to\infty}\int_{-c}^cx\,\diff x=0$ exists while
$\int_{-\infty}^{\infty}x\,\diff x$ does not, and a distribution with such a density
tail has no mean --- the Cauchy distribution being the standard example. The
distinction is exactly the one between the principal value and the integral, and
it is what makes ``the expectation does not exist'' a meaningful statement.

\section{Multivariate Integration}\label{sec:int-multi}

Two results are needed for the Gaussian computation. Both are stated without
proof and used only where cited; complete treatments are in
\textcite[Ch.~10]{rudin1976principles} and, in the measure-theoretic form,
\textcite[Ch.~5]{legall2022measure}.

\begin{theorem}[Fubini--Tonelli]\label{thm:fubini}
	Let $f\colon\R^m\times\R^n\to\R$ be measurable. Suppose either that
	$f\geq0$, or that the iterated integral of $\abs{f}$ is finite. Then
	\[
		\int_{\R^{m+n}}f=\int_{\R^m}\left(\int_{\R^n}f(x,y)\,\diff y\right)\,\diff x
		=\int_{\R^n}\left(\int_{\R^m}f(x,y)\,\diff x\right)\,\diff y,
	\]
	all three integrals existing and being equal.
\end{theorem}

The non-negativity or absolute integrability hypothesis is essential; without
it the two iterated integrals can exist and differ. In economics the theorem is
what licenses interchanging an expectation over states with an integral over
time, or summing a double sum in either order --- routine steps that are valid
under conditions worth stating rather than assuming.

\begin{theorem}[Change of variables]\label{thm:change-of-variables}
	Let $U,V\subseteq\R^n$ be open and $\Phi\colon U\to V$ a
	continuously differentiable bijection with $\det D\Phi(u)\neq0$ for every
	$u\in U$. Then for every integrable $f\colon V\to\R$,
	\begin{equation}\label{eq:change-of-variables}
		\int_{V}f(x)\,\diff x=\int_{U}f\bigl(\Phi(u)\bigr)\,\bigl|\det D\Phi(u)\bigr|\,\diff u .
	\end{equation}
\end{theorem}

The factor $\abs{\det D\Phi}$ is the local volume magnification of $\Phi$, which
is the interpretation of the determinant established in
\autoref{eg:det-jacobian}: $\Phi$ maps an infinitesimal box of volume $\diff u$ to a
parallelepiped of volume $\abs{\det D\Phi(u)}\,\diff u$
\autocite[slide 93]{fukuda2026calculus}. For $n=1$, $\abs{\det D\Phi}=\abs{\varphi'}$
and \eqref{eq:change-of-variables} is \eqref{eq:substitution} with the absolute
value accounting for orientation, which the one-dimensional statement instead
encodes in the order of the limits.

\section{Differentiating Under the Integral Sign}\label{sec:int-leibniz}

\begin{lemma}[Differentiation under the integral]\label{thm:differentiation-under-integral}
	Let $f\colon[a,b]\times[\underline r,\bar r]\to\R$ be such that
	$f(\cdot,r)$ is continuous on $[a,b]$ for each $r$, and
	$\partial f/\partial r$ exists and is continuous on the whole rectangle. Then
	$r\mapsto\int_a^bf(x,r)\,\diff x$ is differentiable on $(\underline r,\bar r)$ with
	\[
		\deriv{}{r}\int_a^bf(x,r)\,\diff x=\int_a^b\frac{\partial f}{\partial r}(x,r)\,\diff x .
	\]
\end{lemma}

\begin{proof}
	Fix $r$ and let $h\neq0$ be small enough that $r+h$ stays in the interval. By
	\autoref{thm:mvt} applied to $s\mapsto f(x,s)$, for each $x$ there is
	$\theta_x\in(0,1)$ with
	$\bigl(f(x,r+h)-f(x,r)\bigr)/h=\frac{\partial f}{\partial r}(x,r+\theta_xh)$.
	Hence
	\[
		\left|\frac{1}{h}\int_a^b\bigl(f(x,r+h)-f(x,r)\bigr)\,\diff x
		-\int_a^b\frac{\partial f}{\partial r}(x,r)\,\diff x\right|
		\leq(b-a)\sup_{x\in[a,b]}
		\left|\frac{\partial f}{\partial r}(x,r+\theta_xh)
		-\frac{\partial f}{\partial r}(x,r)\right| .
	\]
	The rectangle is compact, so $\partial f/\partial r$ is uniformly continuous on
	it by \autoref{thm:heine-cantor}. The two arguments differ by at most
	$\abs{h}$, uniformly in $x$, so the supremum tends to $0$ as $h\to0$.
\end{proof}

\begin{theorem}[Leibniz formula]\label{thm:leibniz}
	Under the hypotheses of \autoref{thm:differentiation-under-integral}, let
	$A,B\colon[\underline r,\bar r]\to[a,b]$ be continuously differentiable and set
	$V(r)\eqdef\int_{A(r)}^{B(r)}f(x,r)\,\diff x$. Then $V$ is differentiable with
	\begin{equation}\label{eq:leibniz}
		V'(r)=f\bigl(B(r),r\bigr)B'(r)-f\bigl(A(r),r\bigr)A'(r)
		+\int_{A(r)}^{B(r)}\frac{\partial f}{\partial r}(x,r)\,\diff x .
	\end{equation}
\end{theorem}

\begin{proof}
	Let $g(x,r)\eqdef\int_a^xf(u,r)\,\diff u$, so that
	\[
		\pderiv{g}{x}=f(x,r)\ \text{ by \autoref{thm:ftc}(i)},
		\qquad
		\pderiv{g}{r}=\int_a^x\frac{\partial f}{\partial r}(u,r)\,\diff u
		\ \text{ by \autoref{thm:differentiation-under-integral}};
	\]
	both are continuous, so $g$ is
	differentiable in the sense of \autoref{def:total-derivative} by
	\autoref{prop:c1-differentiable}. By
	\autoref{prop:integral-properties}(ii), $V(r)=g(B(r),r)-g(A(r),r)$, and the
	chain rule (\autoref{thm:chain-rule}) gives
	\[
		V'(r)=\frac{\partial g}{\partial x}(B(r),r)B'(r)+\frac{\partial g}{\partial r}(B(r),r)
		-\frac{\partial g}{\partial x}(A(r),r)A'(r)-\frac{\partial g}{\partial r}(A(r),r).
	\]
	Substituting the two partial derivatives of $g$ and combining the two
	$\partial g/\partial r$ terms into a single integral from $A(r)$ to $B(r)$
	yields \eqref{eq:leibniz}.
\end{proof}

This is Theorem~3.3 of \textcite{fukuda2026notes}, and the proof is the route
mapped out in Exercise~3.6 there
\autocite[slides 95--96]{fukuda2026calculus}, carried out. The intuitive
derivation on \autocite[slides 97--98]{fukuda2026calculus} identifies the three
terms correctly --- the value gained at the moving upper endpoint, the value lost
at the moving lower endpoint, and the change in the integrand over the fixed
region --- and the proof above supplies the hypotheses under which the
approximations become identities.

\section{Economic Application: Continuation Values}\label{sec:int-continuation}

\begin{proposition}[Differential form of a continuation value]\label{prop:continuation-value}
	Let $f\colon\R\to\R$ be bounded and continuous and $\rho>0$, and define the
	continuation value
	\[
		V(t)\eqdef\int_t^{\infty}e^{-\rho(s-t)}f(s)\,\diff s .
	\]
	Then the integral converges absolutely for every $t$, $V$ is differentiable, and
	\begin{equation}\label{eq:hjb-scalar}
		V'(t)=\rho V(t)-f(t),
		\qquad\text{equivalently}\qquad
		\rho V(t)=f(t)+V'(t).
	\end{equation}
\end{proposition}

\begin{proof}
	Let $\abs{f}\leq M$. Then
	$\int_t^{\infty}e^{-\rho(s-t)}\abs{f(s)}\,\diff s\leq M\int_t^{\infty}e^{-\rho(s-t)}\,\diff s
		=M/\rho<\infty$, so the integral converges absolutely.

	Write $W(t)\eqdef\int_t^{\infty}e^{-\rho s}f(s)\,\diff s$, which converges absolutely
	by the same bound, so that $V(t)=e^{\rho t}W(t)$. Fixing any $t_0$,
	$W(t)=W(t_0)-\int_{t_0}^{t}e^{-\rho s}f(s)\,\diff s$, and the integrand is
	continuous, so \autoref{thm:ftc}(i) gives $W'(t)=-e^{-\rho t}f(t)$. By the
	product rule (\autoref{prop:deriv-product}),
	\[
		V'(t)=\rho e^{\rho t}W(t)+e^{\rho t}W'(t)=\rho V(t)-f(t).
	\]
\end{proof}

\paragraph{Reading \eqref{eq:hjb-scalar}.} In the form
$\rho V=f+V'$ the identity is an asset-pricing equation: the required return
$\rho V$ on holding the ``asset'' whose value is $V$ equals the flow payoff $f$
plus the capital gain $V'$. It is the deterministic, uncontrolled special case of
the Hamilton--Jacobi--Bellman equation of \autoref{ch:continuoustime}, and the
same three terms reappear there with a maximisation over controls inserted in
front of $f$. This is Example~3.22 of \textcite{fukuda2026notes}
\autocite[slide 94]{fukuda2026calculus}.

The proof deliberately avoids differentiating under an improper integral. The
Leibniz formula \eqref{eq:leibniz} applies on compact rectangles, and $[t,\infty)$
is not one; a direct application would require a domination argument. Rewriting
$V=e^{\rho t}W$ moves the $t$-dependence out of the integrand entirely, leaving
only a moving endpoint, and \autoref{thm:ftc}(i) handles that with no additional
hypotheses.

\section{Economic Application: Moments of Two Distributions}\label{sec:int-moments}

\subsection{The exponential distribution}\label{sec:int-exponential}

\begin{proposition}\label{prop:exponential-moments}
	Let $X$ have density $f(x)=\lambda e^{-\lambda x}$ for $x\geq0$ and $f(x)=0$ for
	$x<0$, with $\lambda>0$. Then its distribution function is
	$F(x)=1-e^{-\lambda x}$ for $x>0$, and
	\[
		\E[X]=\frac{1}{\lambda},\qquad \E[X^2]=\frac{2}{\lambda^2},
		\qquad \Var[X]=\frac{1}{\lambda^2} .
	\]
\end{proposition}

\begin{proof}
	For $x>0$, $F(x)=\int_0^x\lambda e^{-\lambda t}\,\diff t
		=\bigl[-e^{-\lambda t}\bigr]_0^x=1-e^{-\lambda x}$ by \autoref{thm:ftc}(ii).

	For the first moment, write $\lambda e^{-\lambda x}=\bigl(-e^{-\lambda
			x}\bigr)'$ and apply \autoref{thm:by-parts} on $[0,c]$:
	\[
		\int_0^{c}x\lambda e^{-\lambda x}\,\diff x
		=\bigl[-xe^{-\lambda x}\bigr]_0^{c}+\int_0^{c}e^{-\lambda x}\,\diff x
		=-ce^{-\lambda c}+\frac{1-e^{-\lambda c}}{\lambda} .
	\]
	Letting $c\to\infty$ and using $ce^{-\lambda c}\to0$ gives $\E[X]=1/\lambda$.

	For the second moment, the same split with $x^2$ in place of $x$ gives
	\[
		\int_0^{c}x^2\lambda e^{-\lambda x}\,\diff x
		=\bigl[-x^2e^{-\lambda x}\bigr]_0^{c}+2\int_0^{c}xe^{-\lambda x}\,\diff x
		=-c^2e^{-\lambda c}+\frac{2}{\lambda}\int_0^{c}x\lambda e^{-\lambda x}\,\diff x,
	\]
	so letting $c\to\infty$ and substituting the first moment yields
	$\E[X^2]=2/\lambda^2$. Finally
	$\Var[X]=\E[X^2]-(\E[X])^2=2/\lambda^2-1/\lambda^2=1/\lambda^2$ by
	\autoref{prop:variance-identity}.
\end{proof}

These are the computations of \autocite[slides 73--77]{fukuda2026calculus}, with
the improper integrals handled as limits of integrals over $[0,c]$ in accordance
with \autoref{def:improper} rather than by evaluating boundary terms at
$\infty$ directly. The exponential distribution is the continuous-time
counterpart of a constant hazard rate: $\Prob(X>s+t\mid X>s)=e^{-\lambda t}$ does
not depend on $s$, which is why it is the standard model of the arrival of a
price-adjustment opportunity in Calvo pricing and of job separation in search
models.

\subsection{The Gaussian integral}\label{sec:int-gaussian}

\begin{proposition}[Gaussian integral]\label{prop:gaussian-integral}
	\[
		\int_{-\infty}^{\infty}e^{-x^2}\,\diff x=\sqrt{\pi},
		\qquad\text{and more generally}\qquad
		\int_{-\infty}^{\infty}e^{-a(x+b)^2}\,\diff x=\sqrt{\frac{\pi}{a}}
	\]
	for every $a>0$ and $b\in\R$. Consequently the function \eqref{eq:normal-density}
	integrates to $1$ and is a probability density.
\end{proposition}

\begin{proof}
	Let $I\eqdef\int_{-\infty}^{\infty}e^{-x^2}\,\diff x$, which converges because
	$e^{-x^2}\leq e^{-\abs{x}}$ for $\abs{x}\geq1$. The integrand of the product
	$I^2$ is non-negative, so \autoref{thm:fubini} applies and
	\[
		I^2=\left(\int_{\R}e^{-x^2}\,\diff x\right)\left(\int_{\R}e^{-y^2}\,\diff y\right)
		=\int_{\R}\int_{\R}e^{-(x^2+y^2)}\,\diff y\,\diff x .
	\]
	Apply \autoref{thm:change-of-variables} with polar coordinates
	$\Phi(r,\theta)=(r\cos\theta,r\sin\theta)$, a continuously differentiable
	bijection from $U=(0,\infty)\times(0,2\pi)$ onto $\R^2$ minus a half-line, a set
	of measure zero which does not affect the integral. Its Jacobian is
	\[
		D\Phi(r,\theta)=\begin{pmatrix}
			\cos\theta & -r\sin\theta \\
			\sin\theta & r\cos\theta
		\end{pmatrix},
		\qquad
		\det D\Phi=r\bigl(\cos^2\theta+\sin^2\theta\bigr)=r>0 .
	\]
	Since $x^2+y^2=r^2$,
	\[
		I^2=\int_0^{2\pi}\!\!\int_0^{\infty}e^{-r^2}\,r\,\diff r\,\diff \theta
		=2\pi\left[-\tfrac12e^{-r^2}\right]_0^{\infty}=2\pi\cdot\tfrac12=\pi .
	\]
	As $I>0$, $I=\sqrt\pi$. For the general case substitute
	$x=y/\sqrt a-b$ in \autoref{thm:substitution}, so that $\diff x=\,\diff y/\sqrt a$ and
	$a(x+b)^2=y^2$, giving $\int e^{-a(x+b)^2}\,\diff x=I/\sqrt a=\sqrt{\pi/a}$. Taking
	$a=1/(2\sigma^2)$ and $b=-\mu$ gives
	$\int_{\R}\exp\bigl(-(x-\mu)^2/(2\sigma^2)\bigr)\,\diff x=\sqrt{2\pi\sigma^2}$, so
	\eqref{eq:normal-density} integrates to $1$.
\end{proof}

This is the derivation of \autocite[slides 89--93]{fukuda2026calculus}, with the
interchange of the two integrals justified by \autoref{thm:fubini} and the polar
substitution by \autoref{thm:change-of-variables} rather than by the differential
mnemonic $\diff y\,\diff x=\abs{\det J}\,\diff r\,\diff \theta$. The trick is that
$e^{-x^2}$ has no elementary antiderivative, so \autoref{thm:ftc}(ii) is
unavailable; squaring converts the problem to two dimensions, where the radial
symmetry of $e^{-(x^2+y^2)}$ makes the extra factor $r$ from the Jacobian
exactly the derivative needed for a substitution.

\subsection{Mean and variance of the normal}\label{sec:int-normal}

\begin{proposition}\label{prop:normal-moments}
	Let $X$ have the density \eqref{eq:normal-density} with $\sigma>0$. Then
	$\E[X]=\mu$ and $\Var[X]=\sigma^2$.
\end{proposition}

\begin{proof}
	\emph{Mean.} Since the density integrates to $1$
	(\autoref{prop:gaussian-integral}),
	\[
		\E[X]=\int_{\R}xf(x)\,\diff x=\int_{\R}(x-\mu)f(x)\,\diff x+\mu .
	\]
	The remaining integral vanishes. Substituting $t=(x-\mu)/\sigma$, so
	$\diff x=\sigma\,\diff t$ by \autoref{thm:substitution}, turns it into
	$\sigma\int_{\R}t\,\phi(t)\,\diff t$ with $\phi(t)=(2\pi)^{-1/2}e^{-t^2/2}$; the
	integrand is odd and absolutely integrable, so each half contributes the
	negative of the other and the integral is $0$. Equivalently, and without
	appealing to symmetry, note that
	\[
		(x-\mu)f(x)=-\sigma^2\deriv{}{x}f(x),
		\qquad\text{so}\qquad
		\int_{-c}^{c}(x-\mu)f(x)\,\diff x=-\sigma^2\bigl[f(x)\bigr]_{-c}^{c}
		\longrightarrow0
	\]
	as $c\to\infty$, since $f(x)\to0$ at both ends.

	\emph{Variance.} Using the same identity $f'(x)=-(x-\mu)\sigma^{-2}f(x)$ and
	integrating by parts (\autoref{thm:by-parts}) on $[-c,c]$,
	\[
		\int_{-c}^{c}(x-\mu)^2f(x)\,\diff x
		=\int_{-c}^{c}(x-\mu)\bigl(-\sigma^2f'(x)\bigr)\,\diff x
		=\bigl[-\sigma^2(x-\mu)f(x)\bigr]_{-c}^{c}+\sigma^2\int_{-c}^{c}f(x)\,\diff x .
	\]
	The boundary term tends to $0$ because $f$ decays faster than any polynomial
	grows, and the remaining integral tends to $1$. Hence
	$\Var[X]=\E[(X-\mu)^2]=\sigma^2$.
\end{proof}

These are the computations of \autocite[slides 85--88]{fukuda2026calculus}. Two
points of technique are worth isolating. First, the identity
$f'=-(x-\mu)\sigma^{-2}f$ is what makes the normal tractable: the density is its
own antiderivative up to the linear factor, so every moment computation is an
integration by parts that reduces the power by one. Second, both boundary terms
are evaluated as limits over $[-c,c]$; writing
$\bigl[\,\cdot\,\bigr]_{-\infty}^{\infty}$ directly, as the slides do,
presupposes the convergence that \autoref{def:improper} requires be checked, and
here it holds because of the Gaussian decay.

\section{Chapter Summary}\label{sec:int-summary}

The Riemann integral is defined by squeezing a bounded function between step
functions; continuity, boundedness with finitely many discontinuities, and
monotonicity each suffice for integrability, by arguments that use uniform
continuity, a covering of the bad set, and telescoping respectively. The
fundamental theorem has two halves with different hypotheses: the indefinite
integral of an integrable function is Lipschitz and differentiates back to the
integrand wherever the integrand is continuous; and any antiderivative evaluates
the definite integral. Integration by parts is the product rule integrated, and
substitution is the chain rule integrated; in several variables the substitution
rule acquires the factor $\abs{\det D\Phi}$, the local volume magnification
established in \autoref{ch:linear}. Improper integrals are limits and require
both tails to converge separately, which is what makes non-existence of a mean a
meaningful statement.

Differentiation under the integral sign holds when the parameter derivative is
continuous on a compact rectangle, by the mean value theorem and uniform
continuity; with moving endpoints it becomes the Leibniz formula
\eqref{eq:leibniz}, whose three terms are the flows at the two endpoints and the
change in the integrand. Applied to a discounted flow it gives
$\rho V=f+V'$, the deterministic Hamilton--Jacobi--Bellman equation.

The exponential distribution's moments follow from repeated integration by parts;
the normal density integrates to $1$ by the polar-coordinate evaluation of the
Gaussian integral, which needs both Fubini's theorem and the multivariate change
of variables; and its mean and variance follow from the identity
$f'=-(x-\mu)\sigma^{-2}f$ together with one integration by parts each.

\begin{exercise}\label{ex:int-1}
	Show that $f(x)=1$ for $x\in\Q\cap[0,1]$ and $f(x)=0$ otherwise is bounded and
	not Riemann integrable, by computing $L(f,P)$ and $U(f,P)$ for an arbitrary
	partition. Then show that it is Lebesgue integrable with integral $0$, and
	identify which hypothesis of \autoref{prop:integrable} fails.
\end{exercise}

\begin{exercise}\label{ex:int-2}
	Compute the moment generating function
	$M(s)=\E[e^{sX}]$ for $X$ normal with mean $\mu$ and variance $\sigma^2$, by
	completing the square in the exponent and applying
	\autoref{prop:gaussian-integral}. Verify $M'(0)=\mu$ and $M''(0)=\mu^2+\sigma^2$,
	stating the hypothesis under which the two differentiations under the integral
	sign are legitimate.
\end{exercise}

\begin{exercise}\label{ex:int-3}
	A firm's value is $V(t)=\int_t^{T(t)}e^{-\rho(s-t)}\pi(s,t)\,\diff s$, where the
	exit date $T$ and the flow profit $\pi$ both depend on $t$. Apply
	\eqref{eq:leibniz} to compute $V'(t)$, identify the three terms economically,
	and show that when $T$ is chosen optimally the term involving $T'(t)$ vanishes.
	Relate the last observation to the envelope theorem of
	\autoref{ch:optimization}.
\end{exercise}
